\SourceLeaf{398}%
\noindent [also nach 12:]
\[
\begin{aligned}
s\frac{df}{du}={}&t'\{a'(Y\zeta-Z\eta)+b'(Z\xi-X\zeta)+c'(X\eta-Y\xi)\}\\
&+\theta'\{\alpha'(Y\zeta-Z\eta)+\beta'(Z\xi-X\zeta)+\gamma'(X\eta-Y\xi)\}\\
&+t''\{a(Y\zeta-Z\eta)+b(Z\xi-X\zeta)+c(X\eta-Y\xi)\}\\
&+\theta''\{\alpha(Y\eta-Z\eta)+\beta(Z\xi-X\zeta)+\gamma(X\eta-Y\xi)\}.
\end{aligned}
\]
[Setzt man]
\[
\frac{\partial\partial x}{\partial t^2}=a^{(2)},\quad
\frac{\partial\partial x}{\partial t\,\partial\theta}=a^{(3)},\quad
\left[\frac{\partial\partial x}{\partial\theta^2}=a^{(4)},\right]
\]
[so ergibt sich]
\[
\begin{aligned}
a'&=a^{(2)}t'+a^{(3)}\theta'\\
\alpha'&=a^{(3)}t'+a^{(4)}\theta'\\
&\text{etc.}
\end{aligned}
\]
[und wegen der Identität $\frac{\partial a}{\partial\theta}=\frac{\partial\alpha}{\partial t}$:]
\begin{equation}\tag*{[15]}
\left\{\begin{aligned}
aa'&=\tfrac12\frac{\partial(aa)}{\partial t}.t'+\tfrac12\frac{\partial(aa)}{\partial\theta}.\theta'\\
a\alpha'&=\tfrac12\frac{\partial(aa)}{\partial\theta}.t'+\frac{\partial(a\alpha)}{\partial\theta}.\theta'-\tfrac12\frac{\partial(\alpha\alpha)}{\partial t}.\theta'\\
\alpha a'&=-\tfrac12\frac{\partial(aa)}{\partial\theta}.t'+\frac{\partial(a\alpha)}{\partial t}.t'+\tfrac12\frac{\partial(\alpha\alpha)}{\partial t}.\theta'\\
\alpha\alpha'&=\tfrac12\frac{\partial(\alpha\alpha)}{\partial t}.t'+\tfrac12\frac{\partial(\alpha\alpha)}{\partial\theta}.\theta'.
\end{aligned}\right.
\end{equation}
[Ferner galt]
\[
\begin{aligned}s\xi&=at'+\alpha\theta'\\s\eta&=bt'+\beta\theta'\\s\zeta&=ct'+\gamma\theta'\end{aligned}
\]
[und]
\[
\begin{aligned}
\sigma X&=b\gamma-c\beta\\
\sigma Y&=c\alpha-a\gamma\\
\sigma Z&=a\beta-b\alpha.
\end{aligned}
\]
[Somit gilt:]
\[
\begin{aligned}
s\sigma(Y\zeta-Z\eta)&=t'(\alpha P-aQ)+\theta'(\alpha Q-aR)\\
s\sigma(Z\xi-X\zeta)&=t'(\beta P-bQ)+\theta'(\beta Q-bR)\\
s\sigma(X\eta-Y\xi)&=t'(\gamma P-cQ)+\theta'(\gamma Q-cR)
\end{aligned}
\]
